\documentclass{article}
\usepackage{amsmath, booktabs, geometry}
\geometry{margin=2.5cm}
\title{Rolling Stock Scheduling -- Problem Formulation}
\author{Gati Shakti Vishwavidyalaya}
\date{}
\begin{document}
\maketitle

\section{Problem Overview}
Assign $N$ train units to $T$ timetabled trips across $S$ stations over $D$ days.
Every trip must be covered, physical continuity must hold, and operational cost minimized.
This matches the structure in Guth Jarkovsky et al. (2026) adapted to our 5-station German network.

\section{Sets and Parameters}
\begin{tabular}{lll}
\toprule
Symbol & Meaning & Code location \\
\midrule
$U$ & Set of train units & fleet.csv, fleet\_lookup in data\_loader.py \\
$T$ & Set of scheduled trips & trips.csv, trips\_lookup in data\_loader.py \\
$S$ & Set of stations & stations.csv \\
$d_{ij}$ & Distance (km) from station $i$ to $j$ & distance\_matrix.csv, dist\_lookup \\
$t_{ij}$ & Travel time (min) from $i$ to $j$ & dist\_lookup[(i,j)][0] \\
$cap_u$ & Passenger capacity of unit $u$ & fleet.csv capacity column \\
$km\_max_u$ & Max km before maintenance & maintenance\_rules.csv, maint\_lookup \\
$dep_u$ & Departure minute of trip $t$ & trips.csv departure\_min column \\
\bottomrule
\end{tabular}

\section{Decision Variables}
Let $x_{u,t} \in \{0, 1\}$ mean unit $u$ is assigned to trip $t$.
In code: chains dict where chains[unit\_id] = list of trip\_ids.

\section{Objective Function}
\textit{Used in: src/utils/cost.py, function compute\_cost()}

Minimize total operational cost:
\[
  \text{Cost} = w_1 \cdot D_{head} + w_2 \cdot N_{units} + w_3 \cdot V_{maint} + w_4 \cdot P_{uncov} + w_5 \cdot V_{time}
\]

Where:
\begin{itemize}
  \item $D_{head}$ = total empty (deadhead) kilometers traveled
  \item $N_{units}$ = number of train units activated
  \item $V_{maint}$ = total km exceeding maintenance threshold
  \item $P_{uncov}$ = total uncovered passenger demand
  \item $V_{time}$ = number of timing constraint violations
  \item Weights: $w_1=1, w_2=150, w_3=5, w_4=2, w_5=1000$ (DEFAULT\_WEIGHTS in cost.py)
\end{itemize}

\section{Constraints}

\subsection{Trip Coverage}
Every trip must have enough capacity:
\[
  \sum_{u \in U} x_{u,t} \cdot cap_u \geq demand_t \quad \forall t \in T
\]
\textit{Checked in: compute\_cost() via capacity\_supplied dict}

\subsection{Continuity Constraint}
Unit $u$ can start trip $t_2$ only after finishing $t_1$ and traveling to the new origin:
\[
  arr_{t_1} + turnaround_{t_1} + t_{dest(t_1),\, orig(t_2)} \leq dep_{t_2}
\]
\textit{Checked in: src/utils/feasibility.py, function is\_feasible\_arc()}

\subsection{Maintenance Constraint}
Total km for unit $u$ must not exceed maintenance threshold:
\[
  \sum_{t \in chain_u} (d_{head,t} + dist_t) \leq km\_max_u
\]
\textit{Tracked in: compute\_cost() via km\_since\_maint accumulator}

\section{Connection to Paper}
This formulation directly mirrors the rolling stock planning problem in
Guth Jarkovsky et al. (2026), where the objective is to minimize empty kilometers
while satisfying trip coverage, continuity, and maintenance constraints.
Their QAOA approach encodes these constraints as a QUBO energy landscape.
Our algorithms (Greedy, SA, GA, NSGA-II) solve the same formulation classically,
establishing baselines before the QUBO quantum phase.

\end{document}
